\chapter{Paneles biológicos empleados}
\label{anexo:paneles}

Este anexo detalla los paneles génicos que utiliza el pipeline como
referencia externa. Es importante distinguir entre dos usos distintos de
estos paneles:

\begin{itemize}
  \item Paneles usados para \emph{definir un eje biológico o una
  categoría} (por ejemplo, ``contenido de pulmón sano residual''
  medido con marcadores canónicos de alvéolo). Estos paneles se declaran
  al pipeline y \textbf{no son un resultado} del análisis, sino una
  herramienta metodológica.
  \item Paneles usados para \emph{validación externa} (por ejemplo, el
  panel IHC clínico de la OMS). Estos paneles se declaran al pipeline
  únicamente en la fase de comprobación y \textbf{sí se contrastan}
  contra la firma para estimar la cobertura de marcadores canónicos.
\end{itemize}

Los paneles se centralizan en \texttt{agentes/utils\_cohortes.py} y en
\texttt{agentes/paso19\_firma\_validada.py}; cualquier cambio en ellos
se propaga a todo el pipeline y a la interfaz.

\section*{Panel IHC clínico (validación externa)}
\addcontentsline{toc}{section}{Panel IHC clínico (validación externa)}

El panel de inmunohistoquímica diagnóstica recomendado por la OMS para
NSCLC se codifica en el pipeline como el diccionario \texttt{IHC\_CLINICA}:

\begin{lstlisting}[style=python, caption={Panel IHC clínico usado como validación externa.}]
IHC_CLINICA = {
    "Escamoso": [
        "KRT5",  "KRT6A", "KRT6B", "KRT13", "KRT14",
        "TP63",  "DSG3",  "DSC3",  "SOX2",  "PKP1",
        "CALML3", "S100A2",
    ],
    "Adenocarcinoma": [
        "NAPSA", "NKX2-1", "SFTPB", "SFTPC", "SFTPA1",
        "SLC34A2", "MUC1", "CEACAM6",
    ],
}
\end{lstlisting}

Esta lista se contrasta con \texttt{FIRMA\_VALIDADA\_COMPLETA.csv}. La
recuperación resultante es 12 de 12 marcadores escamosos y 6 de 8
marcadores de adenocarcinoma (los dos no recuperados son
\gene{SFTPA1} y \gene{SFTPC}, ambas proteínas del surfactante).

\section*{Panel de pulmón sano (definición de eje)}
\addcontentsline{toc}{section}{Panel de pulmón sano (definición de eje)}

Para estimar el contenido residual de pulmón normal dentro de una
muestra tumoral se emplea un panel canónico de marcadores alveolares
que la literatura clásica asocia al pulmón sano:

\begin{lstlisting}[style=python, caption={Panel de marcadores de pulmón sano.}]
MARCADORES_PULMON_NORMAL = {
    "alveolar_tipo_I":  ["AGER", "CLDN18", "SFTA2"],
    "alveolar_tipo_II": ["SFTPA1", "SFTPA2", "SFTPB", "SFTPC", "SFTPD",
                         "NAPSA", "SLC34A2"],
    "adiposo_normal":   ["FABP4", "ADIPOQ"],
    "vascular":         ["VWF", "PECAM1", "CDH5"],
    "otros":            ["WIF1", "AGT"],
}
\end{lstlisting}

El pipeline calcula un score composición sana por muestra tomando el
promedio de expresión normalizada de estos marcadores, y correlaciona
ese score con el output del clasificador. La correlación negativa
observada (véase figura~\ref{fig:composicion}) confirma que la firma no
mide simplemente composición tisular.

Es importante subrayar que estos genes \emph{no} son biomarcadores
identificados por el framework: se utilizan sólo para definir el eje de
composición. La distinción se aclara explícitamente tanto en la
interfaz como en las conclusiones del capítulo~\ref{cap:conclusiones}
para evitar confusión.

\section*{Panel de proliferación (definición de eje)}
\addcontentsline{toc}{section}{Panel de proliferación (definición de eje)}

El eje de actividad proliferativa se define con un panel canónico
compuesto por genes de ciclo celular y proteínas de replicación de ADN:

\begin{lstlisting}[style=python, caption={Panel de proliferación.}]
MARCADORES_TUMOR_GENERICO = {
    "proliferacion": ["MKI67", "TOP2A", "CCNB1", "BIRC5", "AURKA"],
}
\end{lstlisting}

Este panel se corresponde con marcadores universalmente aceptados de
proliferación tumoral y se emplea sólo para caracterizar el segundo eje
biológico de la firma. Como con el panel alveolar, estos genes no son
un resultado del framework sino una herramienta metodológica.

\section*{Cohortes GEO utilizadas}
\addcontentsline{toc}{section}{Cohortes GEO utilizadas}

La lista completa de cohortes descargadas y su estado se recoge en la
tabla~\ref{tab:cohortes} (capítulo~\ref{cap:metodos}). Las cohortes se
declaran en el fichero \texttt{datasets.txt} de la raíz del repositorio,
que contiene una línea por identificador GEO. Modificar este fichero y
reejecutar \texttt{paso1\_descarga.py} incorpora nuevas cohortes al
pipeline; los pasos posteriores las procesan automáticamente sin cambios
de código, siempre que la cohorte cumpla los criterios de inclusión.
